\documentclass[10pt]{article}
\usepackage[margin=0.75in]{geometry}
\usepackage{enumitem}
\usepackage{titlesec}

\titlespacing*{\section}{0pt}{1ex}{0.4ex}
\setlength{\parskip}{0.3ex}
\setlength{\parindent}{0pt}

\pagestyle{empty}

\begin{document}

\begin{center}
{\large \textbf{Resubmission Changes Summary}}\\[0.3em]
{\small Paper: ``Safe and Programmable GPU Resource Management with eBPF''}
\end{center}

\smallskip

This paper was previously submitted to OSDI'26 (Spring cycle) and rejected. Below we summarize the major revisions.

\section{Reframed Core Design Contribution}

The OSDI submission organized the design around three design principles and three corresponding challenges (C1--C3). The revised paper replaces this with a unified execution model based on \emph{asynchronous resource state machines with BPF-directed lifecycle transitions}, with an explicit comparison against existing BPF extensibility systems (sched\_ext, cache\_ext, XDP). The key distinction: GPU policies require many-to-many, asynchronous, cross-device trigger-effect mappings rather than synchronous 1:1 callbacks. This is supported by a new execution model figure (\S4.1) and a concrete MoE expert offloading example showing three cooperating BPF programs across host and device (\S4.1.2).

\section{Added Agent-Driven Evaluation (New RQ2)}

The OSDI submission presented policies and results without describing how they were developed. The revised paper adds RQ2 documenting that all 59 evaluated policies were generated by an AI coding agent performing 974 benchmark runs and 244 policy edits over 20 days (75\% negative results). We document 50 safety events with zero kernel panics or data corruption. Each case study now includes the agent's iterative exploration process: initial failed attempts, diagnostic steps using device-side tracing, and convergence to final policies.

\section{Evaluation Reorganization}

The evaluation was restructured from 3 RQs to 4 RQs. Single-tenant (RQ1) and multi-tenant (now RQ3) retain the same experiments. Agent-driven exploration (RQ2) is entirely new. The programmability and overhead section (now RQ4) retains the existing NVBit comparison, device-side microbenchmarks, and host runtime overhead measurements, and adds a new SIMT verifier correctness evaluation (18 GPU programs tested, 5 unsafe cases correctly rejected) and dynamic attach latency measurement (273\,ms).

\section{Paper Restructuring and Other Changes}

\begin{itemize}[nosep, leftmargin=1.5em]
  \item Challenges (C1--C3) moved from design to motivation; portability (\S7) condensed into implementation.
  \item Design leads with execution model and worked example rather than interface listings.
  \item Reformatted from USENIX to ACM SIGPLAN (ASPLOS); revised abstract.
\end{itemize}

\end{document}
